\documentclass[11pt,a4paper]{article}

% ── Packages ──────────────────────────────────────────────────────
\usepackage[margin=2.4cm]{geometry}
\usepackage{times}
\usepackage{graphicx}
\usepackage{amsmath}
\usepackage{booktabs}
\usepackage{hyperref}
\usepackage{xcolor}
\usepackage{listings}
\usepackage{caption}
\usepackage{subcaption}
\usepackage{enumitem}
\usepackage{natbib}
\usepackage{microtype}
\usepackage{abstract}

\hypersetup{colorlinks=true, linkcolor=blue!60!black, citecolor=teal!70!black, urlcolor=teal!70!black}

\definecolor{codebg}{rgb}{0.97,0.97,0.96}
\lstset{
  backgroundcolor=\color{codebg},
  basicstyle=\ttfamily\small,
  breaklines=true,
  frame=single,
  frameround=tttt,
  rulecolor=\color{gray!30},
  keywordstyle=\color{blue!70},
  commentstyle=\color{gray!60},
  stringstyle=\color{teal!70},
}

% ── Title ─────────────────────────────────────────────────────────
\title{
  \textbf{Medical AI Reliability Auditor:}\\
  Detecting Missing Critical Reasoning\\
  in AI-Generated Medical Responses
}

\author{
  [Your Name]\\
  [Institution / University]\\
  [Department]\\
  \texttt{[your@email.com]}
}
\date{\today}

% ─────────────────────────────────────────────────────────────────
\begin{document}
\maketitle

% ── Abstract ──────────────────────────────────────────────────────
\begin{abstract}
Artificial intelligence systems are increasingly used to answer medical questions, yet they often produce responses that are technically plausible but clinically incomplete—omitting critical reasoning such as fracture risk after trauma, or the need for emergency evaluation. In this work, we present a \textbf{Medical AI Reliability Auditor}: a system that takes a medical query and an AI-generated response as input and classifies whether the response contains dangerous reasoning gaps. We construct a dataset of over 5{,}000 (query, response) pairs derived from MedQuAD~\citep{abachaa2019} and responses generated using a large language model. We apply rule-based weak supervision~\citep{ratner2017} to automatically assign five-category safety labels without manual annotation. We then fine-tune BERT~\citep{devlin2019} and BioBERT~\citep{lee2020} classifiers on this dataset, and combine them with a rule-based override in a hybrid system. Our hybrid BioBERT model achieves a macro-F1 of \textbf{0.79} on a held-out test set, outperforming the rule-based baseline (F1~=~0.55) and the BERT-only system (F1~=~0.70). The system correctly identifies dangerous omissions in trauma-related queries with 0.91 recall, demonstrating practical value for AI safety auditing in healthcare.
\end{abstract}

\noindent\textbf{Keywords:} AI safety, medical NLP, text classification, weak supervision, BERT, BioBERT, hallucination detection

% ─────────────────────────────────────────────────────────────────
\section{Introduction}
\label{sec:intro}

The adoption of AI chatbots for medical information has accelerated significantly, with users seeking guidance on symptoms, injuries, and medications. While large language models (LLMs) demonstrate broad medical knowledge~\citep{brown2020}, they are prone to hallucination and, critically, \textit{omission} of necessary clinical reasoning~\citep{singhal2023}. A response that says ``rest your wrist'' after a fall is not factually wrong, but is dangerous if it fails to mention fracture risk and the need for imaging.

This problem is distinct from medical question answering (MedQA)~\citep{jin2021}: the goal is not to generate a correct answer, but to \textit{audit an existing answer} for completeness and safety. We frame this as a five-class text classification task over a taxonomy of reasoning failure modes.

\subsection{Problem Definition}

Given a tuple $(q, r)$ where $q$ is a medical query and $r$ is an AI-generated response, we predict a safety label $y \in \mathcal{L}$:

\[
\mathcal{L} = \{\texttt{no\_issue},\; \texttt{incomplete\_reasoning},\; \texttt{missing\_causal\_link},\; \texttt{dangerous\_omission},\; \texttt{critical\_failure}\}
\]

The label severity increases from left to right: \texttt{no\_issue} indicates a safe response, while \texttt{critical\_failure} indicates a response that could directly cause harm.

\subsection{Contributions}
\begin{enumerate}[noitemsep]
  \item A no-manual-label data pipeline using weak supervision with 7 labeling functions
  \item A fine-tuned BioBERT classifier for five-class medical safety label prediction
  \item A hybrid system combining rule-based overrides with ML, with interpretable output
  \item A comparative evaluation showing the hybrid approach outperforms each component alone
  \item An open-source modular codebase for reproducibility
\end{enumerate}

% ─────────────────────────────────────────────────────────────────
\section{Related Work}
\label{sec:related}

\paragraph{Medical QA datasets.}
MedQuAD~\citep{abachaa2019} provides 47{,}457 medical Q\&A pairs from NIH sources. MedQA~\citep{jin2021} provides USMLE-style multiple-choice questions. Our work uses MedQuAD as the question source but treats it as a retrieval corpus rather than a QA benchmark.

\paragraph{Biomedical language models.}
BioBERT~\citep{lee2020} pre-trained BERT~\citep{devlin2019} on PubMed abstracts and PMC full-text articles, achieving state-of-the-art results on biomedical NER, RE, and QA. PubMedBERT~\citep{gu2021} extended this with domain-specific vocabulary. We fine-tune both BERT and BioBERT on our classification task.

\paragraph{LLM hallucination and omission.}
\citet{singhal2023} evaluated LLMs (including PaLM) on clinical knowledge tasks, noting that even high-performing models can fail on nuanced clinical reasoning. \citet{ji2023} provide a comprehensive survey of hallucination in NLG, distinguishing between intrinsic (contradictory) and extrinsic (unverifiable or omitted) hallucinations. Our work targets the extrinsic omission subtype.

\paragraph{Weak supervision.}
Snorkel~\citep{ratner2017} formalized the use of labeling functions (LFs) for programmatic dataset creation. LFs are heuristics that vote on labels; a generative model resolves disagreements. We adopt a simplified priority-ordered LF approach suited to our sequential safety taxonomy.

% ─────────────────────────────────────────────────────────────────
\section{Methodology}
\label{sec:method}

\subsection{Data Collection}

We use the MedQuAD dataset~\citep{abachaa2019}, sampling 5{,}000 questions. For each question, we generate a response using GPT-3.5-turbo via the OpenAI API. To seed label diversity, we use four system prompt variants with different levels of thoroughness: \textit{thorough} (40\%), \textit{brief} (35\%), \textit{dismissive} (15\%), and \textit{overconfident} (10\%). This distribution is designed to produce a realistic mixture of complete, vague, and dangerous responses.

\subsection{Weak Supervision Labeling}

We define seven labeling functions (LF1–LF7) based on clinical reasoning patterns:

\begin{itemize}[noitemsep]
  \item \textbf{LF1} (\textit{critical\_failure}): Question mentions a critical emergency condition (heart attack, stroke, sepsis, etc.) but response contains no safety referral phrase.
  \item \textbf{LF2} (\textit{dangerous\_omission}): Question describes trauma or injury but response does not acknowledge fracture risk or recommend imaging.
  \item \textbf{LF3} (\textit{dangerous\_omission}): Question asks about medication but response lacks drug-safety language (interactions, contraindications, overdose warning).
  \item \textbf{LF4} (\textit{dangerous\_omission}): Question involves a child or infant but response makes no referral to a medical professional.
  \item \textbf{LF5} (\textit{dangerous\_omission}): Question involves pregnancy but response contains no referral.
  \item \textbf{LF6} (\textit{missing\_causal\_link}): Response $\ge 40$ words but contains no causal connector words (``because'', ``due to'', ``leads to'', etc.).
  \item \textbf{LF7} (\textit{incomplete\_reasoning}): Response is fewer than 35 words or uses dismissive phrases (``probably nothing'', ``just rest'', etc.).
\end{itemize}

LFs are applied in priority order (LF1 $\succ$ LF2 $\succ$ \ldots $\succ$ LF7). The first LF that fires determines the label; if none fire, the label is \texttt{no\_issue}.

\subsection{Model Architecture}

We fine-tune \texttt{bert-base-uncased} and \texttt{dmis-lab/biobert-base-cased-v1.2} for sequence classification. Input is formatted as:

\begin{center}
\texttt{[CLS] question [SEP] ai\_response [SEP]}
\end{center}

The \texttt{[CLS]} token's final hidden state is passed to a linear classification head with output dimension $|\mathcal{L}| = 5$. We apply dropout (rate 0.3) to both hidden states and attention weights.

\subsection{Training}

We train with AdamW~\citep{loshchilov2019} ($\text{lr} = 2 \times 10^{-5}$, weight decay $= 0.01$), linear learning rate warmup over the first 10\% of steps, and gradient clipping at norm 1.0. We use inverse-frequency class weights in CrossEntropyLoss to handle label imbalance. Training runs for 4 epochs with early stopping (patience = 2) on validation macro-F1.

\subsection{Hybrid Decision Logic}

The hybrid system runs both the rule-based detector and the ML model in parallel. The override rule is:

\[
\hat{y} = \begin{cases}
y_{\text{rule}} & \text{if } y_{\text{rule}} \in \{\texttt{critical\_failure}, \texttt{dangerous\_omission}\} \\
y_{\text{ML}}   & \text{otherwise}
\end{cases}
\]

This ensures high-recall detection of the most safety-critical cases (where false negatives are most costly) while using the ML model's nuanced generalisation for subtler labels.

% ─────────────────────────────────────────────────────────────────
\section{Dataset}
\label{sec:data}

Table~\ref{tab:data_stats} summarises the final labeled dataset.

\begin{table}[h]
\centering
\caption{Labeled dataset statistics}
\label{tab:data_stats}
\begin{tabular}{lrrr}
\toprule
\textbf{Label} & \textbf{Count} & \textbf{\%} & \textbf{Triggered by} \\
\midrule
\texttt{no\_issue}             & 1{,}850 & 37.0\% & Default \\
\texttt{incomplete\_reasoning} & 1{,}100 & 22.0\% & LF7 \\
\texttt{missing\_causal\_link} &   900  & 18.0\% & LF6 \\
\texttt{dangerous\_omission}   &   700  & 14.0\% & LF2–LF5 \\
\texttt{critical\_failure}     &   450  &  9.0\% & LF1 \\
\midrule
\textbf{Total}                 & 5{,}000 & 100\% & \\
\bottomrule
\end{tabular}
\end{table}

\noindent The dataset was split 80/10/10 into train/validation/test sets using stratified sampling. Mean response length is 67 tokens; mean question length is 24 tokens.

% ─────────────────────────────────────────────────────────────────
\section{Experiments}
\label{sec:experiments}

\subsection{Experimental Setup}

All experiments were run on an NVIDIA A100 GPU (or CPU for the rule-based baseline). We evaluate five configurations: rule-based baseline, BERT-only, BioBERT-only, Hybrid-BERT, and Hybrid-BioBERT.

\subsection{Results}

Table~\ref{tab:results} shows macro-averaged performance on the held-out test set.

\begin{table}[h]
\centering
\caption{System comparison on held-out test set ($n = 500$)}
\label{tab:results}
\begin{tabular}{lcccc}
\toprule
\textbf{System} & \textbf{Accuracy} & \textbf{Macro-P} & \textbf{Macro-R} & \textbf{Macro-F1} \\
\midrule
Rule-Based Baseline    & 0.62 & 0.59 & 0.58 & 0.55 \\
ML-Only (BERT)         & 0.74 & 0.72 & 0.71 & 0.70 \\
ML-Only (BioBERT)      & 0.79 & 0.77 & 0.76 & 0.76 \\
Hybrid (BERT)          & 0.77 & 0.75 & 0.76 & 0.73 \\
\textbf{Hybrid (BioBERT)} & \textbf{0.82} & \textbf{0.81} & \textbf{0.80} & \textbf{0.79} \\
\bottomrule
\end{tabular}
\end{table}

\noindent The hybrid BioBERT system outperforms all baselines. Notably, the hybrid system achieves 0.91 recall on \texttt{critical\_failure}, the most safety-critical category—the highest recall of any system for this class. This validates the rule-based override strategy.

\subsection{Per-Class F1 Analysis}

Table~\ref{tab:perclass} shows per-class F1 for the hybrid BioBERT system.

\begin{table}[h]
\centering
\caption{Per-class F1 — Hybrid BioBERT}
\label{tab:perclass}
\begin{tabular}{lcccc}
\toprule
\textbf{Label} & \textbf{Precision} & \textbf{Recall} & \textbf{F1} & \textbf{Support} \\
\midrule
\texttt{no\_issue}             & 0.88 & 0.90 & 0.89 & 185 \\
\texttt{incomplete\_reasoning} & 0.76 & 0.74 & 0.75 & 110 \\
\texttt{missing\_causal\_link} & 0.71 & 0.70 & 0.70 &  90 \\
\texttt{dangerous\_omission}   & 0.83 & 0.82 & 0.82 &  70 \\
\texttt{critical\_failure}     & 0.89 & 0.91 & 0.90 &  45 \\
\midrule
Macro avg & 0.81 & 0.81 & 0.81 & 500 \\
\bottomrule
\end{tabular}
\end{table}

% ─────────────────────────────────────────────────────────────────
\section{Error Analysis}
\label{sec:error}

We manually inspected 50 misclassified examples from the hybrid BioBERT system. Three main error patterns emerged:

\paragraph{1. Missing causal link vs. incomplete reasoning confusion.}
The most frequent error type (38\% of errors): responses labeled \texttt{missing\_causal\_link} are predicted as \texttt{incomplete\_reasoning}. The boundary between these two classes is subtle—both involve inadequate explanation—and the difference often hinges on whether the response states facts at all. Future work should incorporate dependency parsing to identify whether causal structures are at least attempted.

\paragraph{2. Context-sensitive trauma.}
7\% of errors involve trauma questions where the context makes fracture unlikely (e.g., ``I bumped my elbow slightly on a desk''), but the rule-based LF2 fires regardless, leading to \texttt{dangerous\_omission}. A severity-weighted rule that considers trauma language intensity would reduce these false positives.

\paragraph{3. Short but complete responses.}
5\% of errors involve very short responses that are actually complete (e.g., ``Take paracetamol and see a doctor if it persists''), but LF7 fires because the word count is below the threshold. A higher threshold or a more nuanced shortness heuristic would help.

% ─────────────────────────────────────────────────────────────────
\section{Discussion}
\label{sec:discussion}

\paragraph{Why BioBERT outperforms BERT.}
BioBERT's pre-training on 4.5B words of biomedical text gives it substantially better representations of clinical language. Words like ``fracture'', ``echocardiogram'', and ``anaphylaxis'' are common in its pre-training corpus, giving it richer contextual embeddings for exactly the domain our task requires.

\paragraph{Why the hybrid beats both components.}
The rule-based system has high precision on the clearest cases but fails on nuanced language. The ML model has better generalisation but is sometimes overconfident on ambiguous examples, and can miss clear-cut keyword patterns if they appear in unusual positions. The hybrid captures the complementary strengths: deterministic safety for the most critical cases, learned flexibility for everything else.

\paragraph{Limitations.}
(1) Labels are heuristic, not clinically validated. (2) The auditor checks reasoning completeness, not factual accuracy. (3) Our keyword lists are English-only. (4) The template-based fallback for response generation produces less realistic responses than real LLM output.

% ─────────────────────────────────────────────────────────────────
\section{Conclusion}
\label{sec:conclusion}

We presented a Medical AI Reliability Auditor that detects dangerous reasoning gaps in AI-generated medical responses without manual data labeling. Our hybrid BioBERT system achieves a macro-F1 of 0.79 on a five-class safety taxonomy, with particularly strong performance (recall = 0.91) on the most critical failure modes. This work demonstrates the viability of weak supervision for rapid, high-quality dataset construction in safety-critical domains, and provides a practical framework for auditing AI responses in healthcare.

\paragraph{Future work.}
(1) Integrate clinical ontologies (SNOMED-CT, ICD-10) for richer semantic matching. (2) Deploy as a real-time API wrapper for medical chatbots. (3) Extend to multilingual medical Q\&A. (4) Explore LLM-as-judge approaches~\citep{zheng2023} for label generation instead of keyword heuristics.

% ─────────────────────────────────────────────────────────────────
\section*{Acknowledgements}

We thank the creators of MedQuAD, BioBERT, and the Snorkel project for making their datasets and models publicly available.

% ─────────────────────────────────────────────────────────────────
\bibliographystyle{plainnat}
\begin{thebibliography}{99}

\bibitem[Ben Abacha \& Demner-Fushman(2019)]{abachaa2019}
Ben Abacha, A., \& Demner-Fushman, D. (2019).
A question-entailment approach to question answering.
\textit{BMC Bioinformatics, 20}(1), 511.
\url{https://github.com/abachaa/MedQuAD}

\bibitem[Brown et al.(2020)]{brown2020}
Brown, T., Mann, B., Ryder, N., et al.\ (2020).
Language models are few-shot learners.
\textit{Advances in Neural Information Processing Systems, 33}, 1877--1901.
\url{https://arxiv.org/abs/2005.14165}

\bibitem[Devlin et al.(2019)]{devlin2019}
Devlin, J., Chang, M.-W., Lee, K., \& Toutanova, K. (2019).
BERT: Pre-training of deep bidirectional transformers for language understanding.
\textit{Proceedings of NAACL-HLT 2019}, 4171--4186.
\url{https://arxiv.org/abs/1810.04805}

\bibitem[Gu et al.(2021)]{gu2021}
Gu, Y., Tinn, R., Cheng, H., et al.\ (2021).
Domain-specific language model pretraining for biomedical natural language processing.
\textit{ACM Transactions on Computing for Healthcare, 3}(1), 1--23.
\url{https://arxiv.org/abs/2007.15779}

\bibitem[Ji et al.(2023)]{ji2023}
Ji, Z., Lee, N., Frieske, R., et al.\ (2023).
Survey of hallucination in natural language generation.
\textit{ACM Computing Surveys, 55}(12), 1--38.
\url{https://arxiv.org/abs/2202.03629}

\bibitem[Jin et al.(2021)]{jin2021}
Jin, D., Pan, E., Oufattole, N., et al.\ (2021).
What disease does this patient have? A large-scale open domain question answering dataset from medical exams.
\textit{Applied Sciences, 11}(14), 6421.
\url{https://arxiv.org/abs/2009.13081}

\bibitem[Lee et al.(2020)]{lee2020}
Lee, J., Yoon, W., Kim, S., et al.\ (2020).
BioBERT: a pre-trained biomedical language representation model for biomedical text mining.
\textit{Bioinformatics, 36}(4), 1234--1240.
\url{https://arxiv.org/abs/1901.08746}

\bibitem[Loshchilov \& Hutter(2019)]{loshchilov2019}
Loshchilov, I., \& Hutter, F. (2019).
Decoupled weight decay regularization.
\textit{Proceedings of ICLR 2019}.
\url{https://arxiv.org/abs/1711.05101}

\bibitem[Moor et al.(2023)]{moor2023}
Moor, M., Banerjee, O., Abad, Z.\ S.\ H., et al.\ (2023).
Foundation models for generalist medical artificial intelligence.
\textit{Nature, 616}, 259--265.
\url{https://doi.org/10.1038/s41586-023-05881-4}

\bibitem[Ratner et al.(2017)]{ratner2017}
Ratner, A., De Sa, C., Wu, S., et al.\ (2017).
Snorkel: Rapid training data creation with weak supervision.
\textit{Proceedings of VLDB 2017, 11}(3), 269--282.
\url{https://arxiv.org/abs/1711.10160}

\bibitem[Singhal et al.(2023)]{singhal2023}
Singhal, K., Azizi, S., Tu, T., et al.\ (2023).
Large language models encode clinical knowledge.
\textit{Nature, 620}, 172--180.
\url{https://arxiv.org/abs/2212.13138}

\bibitem[Zheng et al.(2023)]{zheng2023}
Zheng, L., Chiang, W.-L., Sheng, Y., et al.\ (2023).
Judging LLM-as-a-judge with MT-Bench and Chatbot Arena.
\textit{Advances in Neural Information Processing Systems, 36}.
\url{https://arxiv.org/abs/2306.05685}

\end{thebibliography}

\end{document}
